% !TEX spellcheck = en_US



% ~~~~~~~~~~~~~~~~~~~~~~~~~~~~~~~~~~~~~~~ V
% (NC) 2026 Haluk Bingol
% github.com/halukbingol/LaTeX-Templates
%
% Licensees may copy, distribute, display, and perform the work and make derivative works 
% and remixes based on it only for non-commercial purposes. 
% ~~~~~~~~~~~~~~~~~~~~~~~~~~~~~~~~~~~~~~~ A




% =====================================================================
%  The Eclipse IDE -- Learning by Example
%  ---------------------------------------------------------------------
%  Built on hblib/Template/hblibTemplate.tex with the libraries in ../../hblib/.
%  Programs and their real outputs are in codeoutput/:
%     codeoutput/<Name>.java  a program (classes without main() are only compiled)
%     codeoutput/<Name>.sh    a terminal session, run in a copy of codeoutput/<Name>/
%     codeoutput/<Name>.txt   what it printed, made by:  bash run_examples.sh
% =====================================================================

\documentclass[aspectratio=169,10pt,compress]{beamer}

% ~~~~~~~~~~~~~~~~~~~~~~~~~~~~~~~~~~~~~~~~~~~~~~~~~~~~~~~~~~~~~~~~~ V hblib
	\usepackage{../../hblib/hblibPresentation} % lib presentation: theme, i/N, \hSection, algorithm2e, ...
	\usepackage{../../hblib/hblibCodeoutput} % lib code + output: \lstcode, \lstterminal, ...
	\usepackage{../../hblib/hblibDatastructures} % lib data structures: \hString, \hStack, \hQueue
% ~~~~~~~~~~~~~~~~~~~~~~~~~~~~~~~~~~~~~~~~~~~~~~~~~~~~~~~~~~~~~~~~~ A hblib


% theme, i/N, fonts (T1, lmodern), booktabs, xcolor, array (L{width} columns), listings, tikz,
% algorithm2e (\small captions, \Fn) and the color gitred come from the hblib libraries




% +++++++++++++++++++++++++++++++++++++++ V specific
\DontPrintSemicolon % in algorithms
\newcommand{\win}{\textbf{Windows}}
\newcommand{\mac}{\textbf{macOS}}
% +++++++++++++++++++++++++++++++++++++++ A specific

% ---------------------------------------------------------------------
\title{The Eclipse IDE}
\subtitle{Learning by Example}
\author[Bingol]{Haluk O. Bingol}

\institute[FCIS]{
    Faculty of Computer and Information Sciences\\
    Yeditepe University
}
\date{
	{\tiny \hbTimeStamp}
}

\begin{document}




% ~~~~~~~~~~~~~~~~~~~~~~~~~~~~~~~~~~~~~~ V frm
\begin{frame}[t,fragile]
  \titlepage
\end{frame}
% ~~~~~~~~~~~~~~~~~~~~~~~~~~~~~~~~~~~~~~ A frm




% ~~~~~~~~~~~~~~~~~~~~~~~~~~~~~~~~~~~~~~ V frm
\begin{frame}[t,fragile] \frametitle{Outline}
  \begin{columns}[T]
    \begin{column}{0.48\textwidth}
      \tableofcontents[sections={1-4}, hideallsubsections]
    \end{column}
    \begin{column}{0.48\textwidth}
      \tableofcontents[sections={5-8}, hideallsubsections]
    \end{column}
  \end{columns}
\end{frame}
% ~~~~~~~~~~~~~~~~~~~~~~~~~~~~~~~~~~~~~~ A frm

% ~~~~~~~~~~~~~~~~~~~~~~~~~~~~~~~~~~~~~~ V
\hSection[Workspace]{A Workspace on Disk}
% ~~~~~~~~~~~~~~~~~~~~~~~~~~~~~~~~~~~~~~ A




% ~~~~~~~~~~~~~~~~~~~~~~~~~~~~~~~~~~~~~~ V frm
\begin{frame}[t,fragile] \frametitle{What Eclipse writes to the disk}
  \begin{columns}[T]
    \begin{column}{0.52\textwidth}
      \lstterminal[basicstyle=\ttfamily\scriptsize]{codeoutput/Workspace.txt}
    \end{column}
    \begin{column}{0.44\textwidth}
      \begin{itemize}\small
        \item A workspace with one project \texttt{Hello} after \emph{File > New > Java Project}, a class and one save.
        \item \texttt{.metadata}: the workspace's own data (here shortened to one file).
        \item \texttt{.classpath}: source folder, JRE, output folder.
        \item Delete \texttt{bin/} and Eclipse rebuilds it: never edit it.
      \end{itemize}
    \end{column}
  \end{columns}
\end{frame}
% ~~~~~~~~~~~~~~~~~~~~~~~~~~~~~~~~~~~~~~ A frm

% ~~~~~~~~~~~~~~~~~~~~~~~~~~~~~~~~~~~~~~ V
\hSection[Build]{Why the Builder Recompiles Other Files}
% ~~~~~~~~~~~~~~~~~~~~~~~~~~~~~~~~~~~~~~ A




% ~~~~~~~~~~~~~~~~~~~~~~~~~~~~~~~~~~~~~~ V frm
\begin{frame}[t,fragile] \frametitle{Two classes, one changes}
  \begin{columns}[T]
    \begin{column}{0.48\textwidth}
      \lstcode[basicstyle=\ttfamily\scriptsize]{codeoutput/show/Shop.java}{Java}
      \lstcode[basicstyle=\ttfamily\scriptsize]{codeoutput/show/Buy.java}{Java}
    \end{column}
    \begin{column}{0.48\textwidth}
      \lstcode[basicstyle=\ttfamily\scriptsize]{codeoutput/show/Shop2.java}{Java}
      \small The new \texttt{Shop}: \texttt{price} now returns a \texttt{double}. Its \textbf{public shape} changed.
      What happens if only \texttt{Shop.java} is compiled again?
    \end{column}
  \end{columns}
\end{frame}
% ~~~~~~~~~~~~~~~~~~~~~~~~~~~~~~~~~~~~~~ A frm




% ~~~~~~~~~~~~~~~~~~~~~~~~~~~~~~~~~~~~~~ V frm
\begin{frame}[t,fragile] \frametitle{Compiling only the changed file is not enough}
  \lstterminal[basicstyle=\ttfamily\scriptsize]{codeoutput/Stale.txt}

  \vspace{0.1em}
  \small
  \texttt{Buy.class} still calls \texttt{int price(int)}, which no longer exists. Eclipse's incremental builder
  recompiles every file that uses a changed signature, so this cannot happen inside Eclipse.
\end{frame}
% ~~~~~~~~~~~~~~~~~~~~~~~~~~~~~~~~~~~~~~ A frm

% ~~~~~~~~~~~~~~~~~~~~~~~~~~~~~~~~~~~~~~ V
\hSection[Problems]{The Problems View}
% ~~~~~~~~~~~~~~~~~~~~~~~~~~~~~~~~~~~~~~ A




% ~~~~~~~~~~~~~~~~~~~~~~~~~~~~~~~~~~~~~~ V frm
\begin{frame}[t,fragile] \frametitle{Errors and warnings of one file}
  \begin{columns}[T]
    \begin{column}{0.46\textwidth}
      \lstcode[basicstyle=\ttfamily\scriptsize]{codeoutput/Markers/Temperature.java}{Java}
    \end{column}
    \begin{column}{0.52\textwidth}
      \uncover<2->{\lstterminal[basicstyle=\ttfamily\tiny]{codeoutput/Markers.txt}}
    \end{column}
  \end{columns}
  \small One error (red marker) and two warnings (yellow markers): the same list as Eclipse's Problems view.
\end{frame}
% ~~~~~~~~~~~~~~~~~~~~~~~~~~~~~~~~~~~~~~ A frm

% ~~~~~~~~~~~~~~~~~~~~~~~~~~~~~~~~~~~~~~ V
\hSection[Run]{Run Configurations}
% ~~~~~~~~~~~~~~~~~~~~~~~~~~~~~~~~~~~~~~ A




% ~~~~~~~~~~~~~~~~~~~~~~~~~~~~~~~~~~~~~~ V frm
\begin{frame}[t,fragile] \frametitle{Program arguments and VM arguments}
  \begin{columns}[T]
    \begin{column}{0.5\textwidth}
      \lstcode[basicstyle=\ttfamily\scriptsize]{codeoutput/show/Check.java}{Java}
      \begin{itemize}\small
        \item \emph{Program arguments} (\texttt{20}, \texttt{-5}) go to \texttt{args}.
        \item \emph{VM arguments} (\texttt{-ea}) go to the JVM: \texttt{-ea} turns \texttt{assert} on.
      \end{itemize}
    \end{column}
    \begin{column}{0.46\textwidth}
      \uncover<2->{\lstterminal[basicstyle=\ttfamily\tiny]{codeoutput/RunConfig.txt}}
    \end{column}
  \end{columns}
\end{frame}
% ~~~~~~~~~~~~~~~~~~~~~~~~~~~~~~~~~~~~~~ A frm

% ~~~~~~~~~~~~~~~~~~~~~~~~~~~~~~~~~~~~~~ V
\hSection[Debug]{The Debug View: the Call Stack}
% ~~~~~~~~~~~~~~~~~~~~~~~~~~~~~~~~~~~~~~ A




% ~~~~~~~~~~~~~~~~~~~~~~~~~~~~~~~~~~~~~~ V frm
\begin{frame}[t,fragile] \frametitle{A recursive method}
  \begin{columns}[T]
    \begin{column}{0.4\textwidth}
      \lstcode[basicstyle=\ttfamily\scriptsize]{codeoutput/Stack/Fact.java}{Java}
      \small A breakpoint on line 3, then \texttt{run}, \texttt{where}, \texttt{locals} (commands in \texttt{stack.txt}).
      \texttt{where} lists one line per active call: the \textbf{call stack}, as in Eclipse's \textbf{Debug} view.
    \end{column}
    \begin{column}{0.58\textwidth}
      \uncover<2->{\lstterminal[basicstyle=\ttfamily\tiny, firstline=12, lastline=30]{codeoutput/Stack.txt}}
    \end{column}
  \end{columns}
\end{frame}
% ~~~~~~~~~~~~~~~~~~~~~~~~~~~~~~~~~~~~~~ A frm




% ~~~~~~~~~~~~~~~~~~~~~~~~~~~~~~~~~~~~~~ V frm
\begin{frame}[t,fragile] \frametitle{The deepest call: four frames of \texttt{fact}}
  \begin{columns}[T]
    \begin{column}{0.4\textwidth}
      \lstcode[basicstyle=\ttfamily\scriptsize]{codeoutput/Stack/Fact.java}{Java}
      \small Three more \emph{Resume}s (\texttt{cont}): \texttt{n} is 3, 2, 1. At \texttt{n = 1} the stack has one
      frame per active call. In Eclipse, click a frame to see its \texttt{n}.
    \end{column}
    \begin{column}{0.58\textwidth}
      \lstterminal[basicstyle=\ttfamily\tiny, firstline=30]{codeoutput/Stack.txt}
    \end{column}
  \end{columns}
\end{frame}
% ~~~~~~~~~~~~~~~~~~~~~~~~~~~~~~~~~~~~~~ A frm

% ~~~~~~~~~~~~~~~~~~~~~~~~~~~~~~~~~~~~~~ V
\hSection[Editing]{Quick Fix and Refactoring}
% ~~~~~~~~~~~~~~~~~~~~~~~~~~~~~~~~~~~~~~ A




% ~~~~~~~~~~~~~~~~~~~~~~~~~~~~~~~~~~~~~~ V frm
\begin{frame}[t,fragile] \frametitle{Quick Fix: \emph{Import 'List' (java.util)}}
  \begin{columns}[T]
    \begin{column}{0.48\textwidth}
      \lstcode[basicstyle=\ttfamily\tiny]{codeoutput/QuickFix/before/Names.java}{Java}
      \lstcode[basicstyle=\ttfamily\tiny]{codeoutput/QuickFix/after/Names.java}{Java}
    \end{column}
    \begin{column}{0.5\textwidth}
      \uncover<2->{\lstterminal[basicstyle=\ttfamily\tiny]{codeoutput/QuickFix.txt}}
      \small \keys{Ctrl+1} / \keys{\cmd+1} on the error, or \keys{Ctrl+\shift+O} / \keys{\cmd+\shift+O}: Eclipse adds both imports.
    \end{column}
  \end{columns}
\end{frame}
% ~~~~~~~~~~~~~~~~~~~~~~~~~~~~~~~~~~~~~~ A frm




% ~~~~~~~~~~~~~~~~~~~~~~~~~~~~~~~~~~~~~~ V frm
\begin{frame}[t,fragile] \frametitle{Refactoring: Rename, Extract Constant, Extract Method}
  \begin{columns}[T]
    \begin{column}{0.44\textwidth}
      \lstcode[basicstyle=\ttfamily\tiny]{codeoutput/Rename/before/Circle.java}{Java}
      \lstcode[basicstyle=\ttfamily\tiny]{codeoutput/Rename/after/Circle.java}{Java}
    \end{column}
    \begin{column}{0.54\textwidth}
      \uncover<2->{\lstterminal[basicstyle=\ttfamily\tiny]{codeoutput/Rename.txt}
      \small Same output, better code. The \texttt{diff} shows every change the three refactorings made.}
    \end{column}
  \end{columns}
\end{frame}
% ~~~~~~~~~~~~~~~~~~~~~~~~~~~~~~~~~~~~~~ A frm

% ~~~~~~~~~~~~~~~~~~~~~~~~~~~~~~~~~~~~~~ V
\hSection[Assist]{Content Assist}
% ~~~~~~~~~~~~~~~~~~~~~~~~~~~~~~~~~~~~~~ A




% ~~~~~~~~~~~~~~~~~~~~~~~~~~~~~~~~~~~~~~ V frm
\begin{frame}[t,fragile] \frametitle{What Content Assist knows}
  \begin{columns}[T]
    \begin{column}{0.5\textwidth}
      \begin{itemize}\small
        \item Type \texttt{Math.s} and press \keys{Ctrl+Space}: Eclipse lists the members of \texttt{Math} that start with \texttt{s}.
        \item It reads them from the class files of the JDK, like \texttt{javap} does on the right.
        \item \texttt{String} alone has more than 100 public members: nobody remembers them all.
      \end{itemize}
    \end{column}
    \begin{column}{0.46\textwidth}
      \uncover<2->{\lstterminal[basicstyle=\ttfamily\tiny]{codeoutput/Assist.txt}}
    \end{column}
  \end{columns}
\end{frame}
% ~~~~~~~~~~~~~~~~~~~~~~~~~~~~~~~~~~~~~~ A frm

% ~~~~~~~~~~~~~~~~~~~~~~~~~~~~~~~~~~~~~~ V
\hSection[Tests]{Running Tests}
% ~~~~~~~~~~~~~~~~~~~~~~~~~~~~~~~~~~~~~~ A




% ~~~~~~~~~~~~~~~~~~~~~~~~~~~~~~~~~~~~~~ V frm
\begin{frame}[t,fragile] \frametitle{A green bar for \texttt{fact}}
  \begin{columns}[T]
    \begin{column}{0.46\textwidth}
      \lstcode[basicstyle=\ttfamily\tiny]{codeoutput/Tests/FactTest.java}{Java}
    \end{column}
    \begin{column}{0.52\textwidth}
      \uncover<2->{\lstterminal[basicstyle=\ttfamily\tiny]{codeoutput/Tests.txt}}
      \small \menu{Run As > JUnit Test} shows the same tree in the JUnit view.
    \end{column}
  \end{columns}
\end{frame}
% ~~~~~~~~~~~~~~~~~~~~~~~~~~~~~~~~~~~~~~ A frm

% ~~~~~~~~~~~~~~~~~~~~~~~~~~~~~~~~~~~~~~ V
\hSection[Summary]{Summary}
% ~~~~~~~~~~~~~~~~~~~~~~~~~~~~~~~~~~~~~~ A




% ~~~~~~~~~~~~~~~~~~~~~~~~~~~~~~~~~~~~~~ V frm
\begin{frame}[t,fragile] \frametitle{Which example shows which idea}
  \begin{center}\small
  \begin{tabular}{ll}
    \toprule
    \textbf{idea (concepts deck)}               & \textbf{example} \\
    \midrule
    workspace, project, \texttt{src}/\texttt{bin} & the files on disk \\
    incremental build                           & \texttt{NoSuchMethodError} when a user of a changed method is not rebuilt \\
    markers, Problems view                      & one error, two warnings \\
    launch configuration                        & program arguments, \texttt{-ea} \\
    debugger, call stack                        & \texttt{where} in a recursion \\
    Quick Fix, refactoring                      & missing imports; Rename, Extract Constant, Extract Method \\
    Content Assist                              & the members of \texttt{Math} \\
    unit tests                                  & JUnit tree \\
    \bottomrule
  \end{tabular}
  \end{center}
\end{frame}
% ~~~~~~~~~~~~~~~~~~~~~~~~~~~~~~~~~~~~~~ A frm




% ~~~~~~~~~~~~~~~~~~~~~~~~~~~~~~~~~~~~~~ V frm
\begin{frame}[t,fragile] \frametitle{Exercises}
  \begin{enumerate}\small
    \item Repeat the \texttt{Shop}/\texttt{Buy} example inside Eclipse: change the return type and watch which files get markers.
    \item In \texttt{Temperature}, fix the error, then make every warning disappear with Quick Fix.
    \item Run \texttt{Check} with \texttt{-ea} from a launch configuration; then turn assertions on for every launch
          (\menu{Preferences > Java > Installed JREs > Edit}, \emph{Default VM arguments}).
    \item Debug \texttt{fact(4)}: at the deepest call, click each line of the Debug view and compare the values of \texttt{n}.
    \item Add a test \texttt{fact(21)} to \texttt{FactTest}. Why does it fail? What would you change?
  \end{enumerate}
\end{frame}
% ~~~~~~~~~~~~~~~~~~~~~~~~~~~~~~~~~~~~~~ A frm




% ~~~~~~~~~~~~~~~~~~~~~~~~~~~~~~~~~~~~~~ V frm
\begin{frame}[t,fragile]
	\begin{center}
		\vspace{4cm}
		\hRemark{\Huge Questions?}
	\end{center}
\end{frame}
% ~~~~~~~~~~~~~~~~~~~~~~~~~~~~~~~~~~~~~~ A frm



\end{document}
